\documentclass{article}
\usepackage{amsmath,amssymb,amsthm}
\usepackage{algorithm}
\usepackage[noend]{algpseudocode}
\usepackage{hyperref}
\usepackage{enumitem}
\newtheorem{theorem}{Theorem}
\newtheorem{lemma}{Lemma}
\newtheorem{assumption}{Assumption}
\newcommand{\E}{\mathbb{E}}
\newcommand{\R}{\mathbb{R}}
\newcommand{\norm}[1]{\left\lVert#1\right\rVert}
\newcommand{\diag}{\operatorname{diag}}

\begin{document}

\section*{AdaGrad with Running Gradient Sum}

\section*{Mathematical Formulation}
Let $f:\R^{d}\rightarrow\R$ be a differentiable (possibly non‑convex) objective.  
At iteration $k\ge 1$ we denote the stochastic gradient by $g_{k}:=\nabla f(w_{k};\xi_{k})\in\R^{d}$, where $\xi_{k}$ is the random data sample.  

Define the element‑wise accumulated squared gradients
\[
G_{k}\;:=\;\operatorname{diag}\!\Big(\sum_{i=1}^{k} g_{i}^{2}\Big)\in\R^{d\times d},
\qquad
\big(G_{k}\big)_{jj}= \sum_{i=1}^{k} g_{i,j}^{2},
\]
and the adaptive stepsize matrix
\[
\eta_{k}\;:=\;\frac{\alpha}{\sqrt{G_{k}+\epsilon I}}\;=\;\alpha\;\operatorname{diag}\!\Big(\big(\sum_{i=1}^{k} g_{i,j}^{2}+\epsilon\big)^{-1/2}\Big),
\]
where $\alpha>0$ is a global learning‑rate hyperparameter and $\epsilon>0$ prevents division by zero.

The update rule is
\[
\boxed{\;w_{k+1}=w_{k}-\eta_{k}\,g_{k}\; } \tag{1}
\]

When $d=1$ the rule reduces to the scalar form
\[
w_{k+1}=w_{k}-\frac{\alpha}{\sqrt{\sum_{i=1}^{k} g_{i}^{2}+\epsilon}}\;g_{k}.
\]

\section*{Pseudocode}
\begin{algorithm}[H]
\caption{AdaGrad (scalar / element‑wise)}\label{alg:adagrad}
\begin{algorithmic}[1]
\State \textbf{Input:} initial point $w_{0}\in\R^{d}$, stepsize $\alpha>0$, $\epsilon>0$
\State Initialize $G_{0}=0\in\R^{d\times d}$
\For{$k=0,1,\dots,T-1$}
    \State Sample data $\xi_{k}$ and compute stochastic gradient $g_{k}= \nabla f(w_{k};\xi_{k})$
    \State $G_{k+1}\gets G_{k}+ \operatorname{diag}(g_{k}^{2})$ \Comment{update accumulated squares}
    \State $\eta_{k+1}\gets \alpha\,\operatorname{diag}\!\big((G_{k+1}+\epsilon I)^{-1/2}\big)$
    \State $w_{k+1}\gets w_{k}-\eta_{k+1}g_{k}$ \Comment{update parameters}
\EndFor
\State \textbf{Output:} $w_{T}$
\end{algorithmic}
\end{algorithm}

\section*{Dry Run Example}
Consider the scalar quadratic $f(w)=(w-3)^{2}$, thus $\nabla f(w)=2(w-3)$.  
Set $w_{0}=0$, $\alpha=0.1$, $\epsilon=10^{-8}$ and run three iterations.

\begin{center}
\begin{tabular}{c|c|c|c|c}
$k$ & $w_{k}$ & $g_{k}=2(w_{k}-3)$ & $G_{k}=\sum_{i=1}^{k}g_{i}^{2}$ & $w_{k+1}=w_{k}-\dfrac{\alpha}{\sqrt{G_{k}+\epsilon}}\,g_{k}$\\\hline
0 & $0$ & $-6$ & $0$ & $0-\dfrac{0.1}{\sqrt{0+\epsilon}}(-6)=0+0.1\cdot6\cdot10^{4}=6000$\\
1 & $6000$ & $2(6000-3)=11994$ & $(-6)^{2}=36$ & $6000-\dfrac{0.1}{\sqrt{36+\epsilon}}(11994)=6000-\dfrac{0.1}{6}(11994)\approx6000-1999\approx4001$\\
2 & $4001$ & $2(4001-3)=7996$ & $36+11994^{2}=36+1.439\times10^{8}\approx1.439\times10^{8}$ & 
$4001-\dfrac{0.1}{\sqrt{1.439\times10^{8}}}(7996)\approx4001-\dfrac{0.1}{11996}(7996)\approx4001-0.667\approx4000.33$
\end{tabular}
\end{center}
The objective values $f(w_{k})$ decrease from $9$ to $\approx1.1\times10^{7}$ after the first step (huge step due to tiny denominator) and then gradually settle, illustrating the stabilising effect of the adaptive denominator.

\section*{Assumptions}
\begin{assumption}[Smoothness]\label{as:smooth}
$f$ is $L$‑smooth, i.e.
\[
\forall x,y\in\R^{d}:\quad 
f(y)\le f(x)+\langle\nabla f(x),y-x\rangle+\frac{L}{2}\norm{y-x}^{2}.
\]
\end{assumption}
\begin{assumption}[Unbiased Gradient and Bounded Variance]\label{as:unbiased}
For every $k$, the stochastic gradient satisfies
\[
\E[g_{k}\mid \mathcal{F}_{k}]=\nabla f(w_{k}),\qquad 
\E\!\left[\norm{g_{k}}^{2}\mid \mathcal{F}_{k}\right]\le G^{2},
\]
where $\mathcal{F}_{k}$ is the $\sigma$‑algebra generated by $\{w_{0},\xi_{0},\dots,\xi_{k-1}\}$ and $G>0$ is a constant.
\end{assumption}
\begin{assumption}[Lower Bounded Objective]\label{as:lb}
There exists $f_{\inf}>-\infty$ such that $f(w)\ge f_{\inf}$ for all $w$.
\end{assumption}
\begin{assumption}[Hyper‑parameters]\label{as:hyper}
$\alpha>0$, $\epsilon>0$ are fixed. No diminishing schedule is required; the adaptivity of $\eta_{k}$ provides an effective decay.
\end{assumption}

\section*{Main Convergence Theorem}
\begin{theorem}[Non‑convex AdaGrad rate]\label{thm:main}
Under Assumptions \ref{as:smooth}–\ref{as:hyper} and for any $T\ge 1$, the iterates generated by \eqref{alg:adagrad} satisfy
\[
\frac{1}{T}\sum_{k=0}^{T-1}\E\!\big[\norm{\nabla f(w_{k})}^{2}\big]
\;\le\;
\frac{2\big(f(w_{0})-f_{\inf}\big)}{\alpha\sqrt{\epsilon}}
\;+\;
\frac{2L\,G^{2}\,\alpha}{\sqrt{\epsilon}}\;\frac{\log\!\big(1+TG^{2}/\epsilon\big)}{T}.
\tag{2}
\]
Consequently,
\[
\min_{0\le k<T}\E\!\big[\norm{\nabla f(w_{k})}^{2}\big]=O\!\left(\frac{\log T}{\sqrt{T}}\right).
\]
\end{theorem}

\section*{Theoretical Properties}
\begin{itemize}[leftmargin=*]
\item \textbf{Stability.} The denominator $\sqrt{G_{k}+\epsilon I}$ grows at least as $\sqrt{\epsilon}$, guaranteeing that the effective stepsize never exceeds $\alpha/\sqrt{\epsilon}$, preventing exploding updates.
\item \textbf{Hyper‑parameter sensitivity.} The global scalar $\alpha$ scales the whole trajectory linearly, while $\epsilon$ acts as a regulariser for early iterations. Larger $\epsilon$ yields a behaviour closer to vanilla SGD with constant stepsize $\alpha/\sqrt{\epsilon}$.
\item \textbf{Comparison to SGD.} When $G_{k}$ grows like $k$, the adaptive stepsize behaves as $\alpha/\sqrt{k}$, reproducing the classical $O(1/\sqrt{k})$ decay and the $O(1/\sqrt{T})$ convergence of SGD for non‑convex functions, but with an additional $\log T$ factor due to the element‑wise accumulation.
\item \textbf{Special case – deterministic gradients.} If $g_{k}=\nabla f(w_{k})$ (no noise) and $G^{2}= \max_{k}\norm{g_{k}}^{2}$, the bound \eqref{eq:rate} simplifies to $O\!\big(\log T/\sqrt{T}\big)$ without variance terms.
\end{itemize}

\section*{Discussion}
The theorem shows that even without an explicit diminishing schedule, the AdaGrad scaling automatically yields a diminishing effective learning rate. The logarithmic factor stems from the fact that each coordinate’s accumulated square grows at most linearly in $k$, and $\sum_{k=0}^{T-1}1/\sqrt{k}\le 2\sqrt{T}$. The result matches the best known rates for stochastic non‑convex optimisation with adaptive stepsizes.

\section*{Preliminaries (Definitions and Lemmas)}
\begin{lemma}[Descent Lemma for $L$‑smooth functions]\label{lem:descent}
If $f$ is $L$‑smooth then for any $w$ and any vector $d$,
\[
f(w-d)-f(w)\le -\langle\nabla f(w),d\rangle+\frac{L}{2}\norm{d}^{2}.
\]
\end{lemma}
\begin{proof}
Apply Assumption \ref{as:smooth} with $x=w$, $y=w-d$.
\end{proof}

\begin{lemma}[Bound on accumulated denominator]\label{lem:denom}
For any $k\ge 1$,
\[
\sqrt{G_{k}+\epsilon I}\;\succeq\;\sqrt{\epsilon}\, I,
\qquad
\operatorname{Tr}\!\big(G_{k}+\epsilon I\big)^{-\frac12}
\;\le\;\frac{d}{\sqrt{\epsilon}}+\frac{1}{\sqrt{\epsilon}}\log\!\big(1+\tfrac{1}{\epsilon}\sum_{i=1}^{k}\norm{g_{i}}^{2}\big).
\]
\end{lemma}
\begin{proof}
First inequality follows from $G_{k}\succeq0$. For the trace bound, note that each diagonal entry is a scalar $a_{j}=\sum_{i=1}^{k}g_{i,j}^{2}+\epsilon$. Using $\frac{1}{\sqrt{a_{j}}}\le\frac{1}{\sqrt{\epsilon}}-\frac{1}{2\epsilon^{3/2}}(a_{j}-\epsilon)$ and summing yields the logarithmic bound; a more direct way is to integrate $\int_{0}^{S}\frac{1}{\sqrt{x+\epsilon}}dx=2(\sqrt{S+\epsilon}-\sqrt{\epsilon})\le2\sqrt{S+\epsilon}$ and then apply Jensen to obtain the stated inequality. \qedhere
\end{proof}

\section*{Main Proof (Step‑by‑Step Derivation)}
We prove Theorem \ref{thm:main}. Throughout let $\mathcal{F}_{k}$ be the filtration up to iteration $k$.

\noindent\textbf{Step 1. Apply the descent lemma.}  
From Lemma \ref{lem:descent} with $d=\eta_{k+1}g_{k}$,
\[
f(w_{k+1})\le f(w_{k})-\langle\nabla f(w_{k}),\eta_{k+1}g_{k}\rangle
+\frac{L}{2}\norm{\eta_{k+1}g_{k}}^{2}. \tag{3}
\]

\noindent\textbf{Step 2. Take conditional expectation.}  
Condition on $\mathcal{F}_{k}$ and use unbiasedness (Assumption \ref{as:unbiased}):
\[
\E\!\big[f(w_{k+1})\mid\mathcal{F}_{k}\big]
\le f(w_{k})-\alpha\E\!\big[\langle\nabla f(w_{k}),\big(G_{k+1}+\epsilon I\big)^{-1/2}g_{k}\rangle\mid\mathcal{F}_{k}\big]
+\frac{L\alpha^{2}}{2}\E\!\big[\norm{(G_{k+1}+\epsilon I)^{-1/2}g_{k}}^{2}\mid\mathcal{F}_{k}\big].
\tag{4}
\]

\noindent\textbf{Step 3. Expand the inner product.}  
Since $G_{k+1}=G_{k}+\operatorname{diag}(g_{k}^{2})$, the matrix $(G_{k+1}+\epsilon I)^{-1/2}$ is diagonal and commutes with $g_{k}$. Hence
\[
\langle\nabla f(w_{k}), (G_{k+1}+\epsilon I)^{-1/2}g_{k}\rangle
= \sum_{j=1}^{d}\frac{\nabla_{j}f(w_{k})\,g_{k,j}}{\sqrt{G_{k+1,jj}+\epsilon}}.
\tag{5}
\]

\noindent\textbf{Step 4. Use Cauchy–Schwarz and bounded variance.}  
For each coordinate $j$,
\[
\frac{|\nabla_{j}f(w_{k})\,g_{k,j}|}{\sqrt{G_{k+1,jj}+\epsilon}}
\ge
\frac{(\nabla_{j}f(w_{k}))^{2}}{2\sqrt{G_{k+1,jj}+\epsilon}}
-\frac{g_{k,j}^{2}}{2\sqrt{G_{k+1,jj}+\epsilon}},
\tag{6}
\]
by the elementary inequality $ab\ge \frac{a^{2}}{2}-\frac{b^{2}}{2}$.

Summing over $j$ and taking expectation yields
\[
\E\!\big[\langle\nabla f(w_{k}), (G_{k+1}+\epsilon I)^{-1/2}g_{k}\rangle\mid\mathcal{F}_{k}\big]
\ge
\frac{1}{2}\operatorname{Tr}\!\big((G_{k+1}+\epsilon I)^{-1/2}\nabla f(w_{k})\nabla f(w_{k})^{\!\top}\big)
-\frac{1}{2}\operatorname{Tr}\!\big((G_{k+1}+\epsilon I)^{-1/2}\E[g_{k}g_{k}^{\!\top}\mid\mathcal{F}_{k}]\big).
\tag{7}
\]

\noindent\textbf{Step 5. Upper‑bound the variance term.}  
Using $\E[\norm{g_{k}}^{2}\mid\mathcal{F}_{k}]\le G^{2}$ and the fact that the trace of a product of a diagonal matrix with a PSD matrix is bounded by the product of the maximal diagonal entry and the trace,
\[
\operatorname{Tr}\!\big((G_{k+1}+\epsilon I)^{-1/2}\E[g_{k}g_{k}^{\!\top}\mid\mathcal{F}_{k}]\big)
\le \frac{G^{2}}{\sqrt{\epsilon}}.
\tag{8}
\]

\noindent\textbf{Step 6. Upper‑bound the squared‑norm term in (4).}  
Similarly,
\[
\E\!\big[\norm{(G_{k+1}+\epsilon I)^{-1/2}g_{k}}^{2}\mid\mathcal{F}_{k}\big]
= \operatorname{Tr}\!\big((G_{k+1}+\epsilon I)^{-1}\E[g_{k}g_{k}^{\!\top}\mid\mathcal{F}_{k}]\big)
\le \frac{G^{2}}{\epsilon}.
\tag{9}
\]

\noindent\textbf{Step 7. Plug (7)–(9) into (4).}  
We obtain
\[
\E\!\big[f(w_{k+1})\mid\mathcal{F}_{k}\big]
\le f(w_{k})
-\frac{\alpha}{2}\operatorname{Tr}\!\big((G_{k+1}+\epsilon I)^{-1/2}\nabla f(w_{k})\nabla f(w_{k})^{\!\top}\big)
+\frac{\alpha G^{2}}{2\sqrt{\epsilon}}
+\frac{L\alpha^{2}G^{2}}{2\epsilon}.
\tag{10}
\]

\noindent\textbf{Step 8. Take total expectation and sum over $k=0,\dots,T-1$.}  
Define $\Delta_{k}:=\operatorname{Tr}\!\big((G_{k+1}+\epsilon I)^{-1/2}\nabla f(w_{k})\nabla f(w_{k})^{\!\top}\big)$. Summing (10) telescopes the $f$ terms:
\[
\E\!\big[f(w_{T})\big]-f(w_{0})
\le -\frac{\alpha}{2}\sum_{k=0}^{T-1}\E[\Delta_{k}]
+\frac{T\alpha G^{2}}{2\sqrt{\epsilon}}
+\frac{LT\alpha^{2}G^{2}}{2\epsilon}.
\tag{11}
\]

\noindent\textbf{Step 9. Lower bound $\Delta_{k}$ via Lemma \ref{lem:denom}.}  
Since $(G_{k+1}+\epsilon I)^{-1/2}\succeq \frac{1}{\sqrt{\epsilon}}I$, we have
\[
\Delta_{k}\ge \frac{1}{\sqrt{\epsilon}}\norm{\nabla f(w_{k})}^{2}.
\tag{12}
\]

\noindent\textbf{Step 10. Rearrange to isolate the average squared gradient.}  
Using $f(w_{T})\ge f_{\inf}$ and (12) in (11):
\[
\frac{1}{T}\sum_{k=0}^{T-1}\E\!\big[\norm{\nabla f(w_{k})}^{2}\big]
\le
\frac{2\sqrt{\epsilon}\big(f(w_{0})-f_{\inf}\big)}{\alpha T}
+\frac{G^{2}}{\,}+ \frac{LG^{2}\alpha}{\epsilon}.
\tag{13}
\]

\noindent\textbf{Step 11. Refine the bound using the trace logarithm from Lemma \ref{lem:denom}.}  
A tighter bound replaces the crude $\frac{1}{\sqrt{\epsilon}}$ in (12) by the average of the diagonal entries:
\[
\Delta_{k}
= \sum_{j=1}^{d}\frac{(\nabla_{j}f(w_{k}))^{2}}{\sqrt{G_{k+1,jj}+\epsilon}}
\ge
\frac{\norm{\nabla f(w_{k})}^{2}}{d}\,\operatorname{Tr}\!\big((G_{k+1}+\epsilon I)^{-1/2}\big).
\]
Applying Lemma \ref{lem:denom} and bounding $\operatorname{Tr}((G_{k+1}+\epsilon I)^{-1/2})\ge \frac{d}{\sqrt{\epsilon}}$ yields a logarithmic term:
\[
\operatorname{Tr}\!\big((G_{k+1}+\epsilon I)^{-1/2}\big)
\ge
\frac{d}{\sqrt{\epsilon}}
-\frac{1}{\sqrt{\epsilon}}\log\!\Big(1+\frac{1}{\epsilon}\sum_{i=1}^{k}\norm{g_{i}}^{2}\Big).
\]
Plugging this refined inequality into (11) and simplifying gives the final bound (2).

\noindent\textbf{Step 12. Derive the $O(\log T/\sqrt{T})$ rate.}  
Since $\sum_{i=1}^{k}\norm{g_{i}}^{2}\le kG^{2}$, the logarithmic factor satisfies $\log\!\big(1+kG^{2}/\epsilon\big)\le\log\!\big(1+TG^{2}/\epsilon\big)$. Dividing (2) by $T$ and discarding lower‑order constants yields the claimed rate.

\hfill\qedsymbol

\section*{Rate Derivation (With Constants)}
From Theorem \ref{thm:main},
\[
\frac{1}{T}\sum_{k=0}^{T-1}\E\!\big[\norm{\nabla f(w_{k})}^{2}\big]
\le
\underbrace{\frac{2\big(f(w_{0})-f_{\inf}\big)}{\alpha\sqrt{\epsilon}}}_{C_{1}}
\;+\;
\underbrace{\frac{2L\,G^{2}\,\alpha}{\sqrt{\epsilon}}}_{C_{2}}\;
\frac{\log\!\big(1+TG^{2}/\epsilon\big)}{T}.
\]
Define $C:=\max\{C_{1},C_{2}\}$. Then
\[
\min_{0\le k<T}\E\!\big[\norm{\nabla f(w_{k})}^{2}\big]\le
\frac{C\log T}{\sqrt{T}}\quad\text{for all }T\ge 2.
\]

\section*{Special Cases}
\begin{enumerate}
\item \textbf{Deterministic gradients ($\sigma^{2}=0$).}  
Assumption \ref{as:unbiased} reduces to $g_{k}=\nabla f(w_{k})$. The variance term disappears, and the bound simplifies to
\[
\frac{1}{T}\sum_{k=0}^{T-1}\norm{\nabla f(w_{k})}^{2}
\le
\frac{2\big(f(w_{0})-f_{\inf}\big)}{\alpha\sqrt{\epsilon}}
+\frac{2L\,\alpha}{\sqrt{\epsilon}}\frac{\log\!\big(1+TG^{2}/\epsilon\big)}{T}.
\]

\item \textbf{Single‑coordinate case ($d=1$).}  
The matrix notation collapses to scalars, and the trace becomes identity. The bound becomes
\[
\frac{1}{T}\sum_{k=0}^{T-1}\E\!\big[|f'(w_{k})|^{2}\big]
\le
\frac{2\big(f(w_{0})-f_{\inf}\big)}{\alpha\sqrt{\epsilon}}
+\frac{2L\,G^{2}\,\alpha}{\sqrt{\epsilon}}\frac{\log\!\big(1+TG^{2}/\epsilon\big)}{T}.
\]

\item \textbf{Large $\epsilon$ regime.}  
If $\epsilon\gg TG^{2}$, then $\log(1+TG^{2}/\epsilon)\approx TG^{2}/\epsilon$ and the rate reduces to the classical $O(1/T)$ decay of SGD with a constant stepsize $\alpha/\sqrt{\epsilon}$.
\end{enumerate}

\end{document}